% ----- Page de Garde -----

\begin{center}
    \includegraphics[width=0.35\textwidth]{image/ISEN.jpg}\\
\end{center}
    \rule{\linewidth}{0.5mm}
        
    \begin{center}
        {\Huge \bf Final Project}
    \end{center}
    
    \begin{center}
        {\LARGE \bf The Maximum Edge Weight Clique Problem}
    \end{center}
    \rule{\linewidth}{0.5mm}
    
    \vspace{1\baselineskip}
    
    \begin{center}
    \begin{tikzpicture}[commutative diagrams/every diagram]
        \node (P0) at (90:4cm) {Alexandre \textsc{Auguste}};
        \node (P1) at (90+72:4cm) {Antoine \textsc{Esnault}} ;
        \node (P2) at (90+2*72:4cm) {Valentin \textsc{Herv\'e}};
        \node (P3) at (90+3*72:4cm) {Martin \textsc{Lobel}};
        \node (P4) at (90+4*72:4cm) {Youn \textsc{M\'elois}};
        \path
        (P0) edge node {} (P1)
        (P0) edge node {} (P2)
        (P0) edge node {} (P3)
        (P0) edge node {} (P4)
        (P1) edge node {} (P2)
        (P1) edge node {} (P3)
        (P1) edge node {} (P4)
        (P2) edge node {} (P3)
        (P2) edge node {} (P4)
        (P3) edge node {} (P4);
    \end{tikzpicture}
    \end{center}
    \vspace{1\baselineskip}   
    \large \textbf{Abstract} The Maximum Edge Weight Clique (MEWC) problem is an optimization problem in graph theory that asks for the clique (a subset of vertices, all adjacent to one another) with the maximum total weight in an edge-weighted undirected graph. In the MEWC problem, each edge has a weight, and the weight of a clique is the sum of the weights of its edges. The goal is to find a clique with the maximum possible weight. 

    \begin{center} \Large
        \emph{Class group :}\\
        \textsc{CIR3} - Team 1 \\
    \end{center}
    
    \begin{center} \Large
        \emph{Teacher :}\\
        Leandro \textsc{Montero}
    \end{center}
    
    \vfill
    
\begin{center} \Large
        \today
\end{center}
    
\newpage